% Algebra 1 - one page of key exercise questions, with a solutions appendix
% Source: Text & Tests 4 (New Edition), Chapter 1 "Algebra 1"
% Compile with: pdflatex algebra1_key_questions.tex
\documentclass[11pt]{article}
\usepackage[a4paper,top=1.5cm,bottom=1.5cm,left=1.6cm,right=1.6cm]{geometry}
\usepackage[T1]{fontenc}
\usepackage{amsmath,amssymb}
\usepackage{textcomp}
\usepackage{xcolor}
\usepackage{soul}
\usepackage{tcolorbox}
\usepackage{enumitem}
\usepackage{titlesec}
\usepackage{multicol}
\usepackage{array,booktabs}
\usepackage{fancyhdr}
\usepackage[hidelinks]{hyperref}

\definecolor{lilac}{RGB}{228,214,248}      % light purple highlight
\definecolor{lilacdeep}{RGB}{104,58,160}
\definecolor{lilacline}{RGB}{176,150,214}
\definecolor{greyish}{RGB}{95,95,105}
\sethlcolor{lilac}

\newcommand{\key}[1]{\hl{#1}}
\newcommand{\keyin}[1]{\textbf{\color{lilacdeep}#1}}  % emphasis inside a purple box
% NB: maths inside \key{...} must be written literally as \mbox{$...$} (soul limitation)

\newtcolorbox{formulabox}[1][]{colback=lilac,colframe=lilacline,boxrule=0.6pt,
  arc=2pt,left=6pt,right=6pt,top=3pt,bottom=3pt,#1}

\titleformat{\section}{\color{lilacdeep}\normalsize\bfseries}{}{0em}{}
\titlespacing*{\section}{0pt}{0.6em}{0.2em}

\pagestyle{fancy}\fancyhf{}
\renewcommand{\headrulewidth}{0pt}
\fancyfoot[L]{\footnotesize\color{greyish}Algebra 1 --- key questions and solutions}
\fancyfoot[R]{\footnotesize\color{greyish}\thepage}

\setlength{\parindent}{0pt}
\setlength{\parskip}{0.3em}

% question list: tight, numbered Q1, Q2, ...
\newlist{qlist}{enumerate}{1}
\setlist[qlist]{label=\textbf{\color{lilacdeep}Q\arabic*.},leftmargin=1.9em,
  topsep=1pt,itemsep=1pt,parsep=0pt}
\newcommand{\bk}[1]{{\footnotesize\color{greyish}[#1]}}

\newcommand{\solhead}[2]{\medskip{\color{lilacdeep}\bfseries Q#1.}\ \ {\color{greyish}\footnotesize #2}\par\nobreak}

\begin{document}

%%%%%%%%%%%%%%%%%%%%%%%%%%%%%%%%  PAGE 1: QUESTIONS  %%%%%%%%%%%%%%%%%%%%%%%%%%%%%%%%

{\color{lilacdeep}\large\bfseries Algebra 1 --- Thirty-Three Key Questions}\hfill
{\color{greyish}\footnotesize \textit{Text \& Tests 4} (New Edition), Ch.\ 1 $\cdot$ LC Higher Level}\\[-5pt]

\textcolor{lilacline}{\rule{\textwidth}{0.9pt}}\\[-6pt]

{\footnotesize One question for each skill the chapter tests, taken from Exercises 1.1--1.10 and the
Revision Exercises; the reference after each is where it comes from. \key{Full solutions in the appendix overleaf.}}

\begin{multicols}{2}
\small
{\color{lilacdeep}\footnotesize\bfseries POLYNOMIAL EXPRESSIONS (1.1)}
\begin{qlist}[font=\bfseries\color{lilacdeep}]
\item Give two reasons why $3x^{2}-\dfrac4x+x^{\frac32}$ is not a polynomial. \bk{Ex 1.1 Q3}
\item If $25x^{2}+tx+4$ is a perfect square for all $x$, find $t$. \bk{Ex 1.1 Q11}
\item Find the coefficient of $x^{2}$ in $(3x^{2}+5x-1)(2x^{2}-6x-5)$. \bk{Ex 1.1 Q19}
\item Divide $x^{3}-8x^{2}+19x-12$ by $(x-1)$. \bk{Ex 1.1 Q24}
\item Divide $8x^{3}-27y^{3}$ by $(2x-3y)$. \bk{Ex 1.1 Q26}
\end{qlist}

{\color{lilacdeep}\footnotesize\bfseries POLYNOMIAL FUNCTIONS (1.2)}
\begin{qlist}[resume]
\item The number of diagonals of an $n$-sided polygon is $d(n)=\dfrac{n^{2}}{2}-\dfrac{3n}{2}$. Find $d(4)$, $d(5)$, $d(6)$, and explain why $d(3)=0$. \bk{Ex 1.2 Q8}
\item Given $f(x)=x^{2}-3x+6$, find $f(-2t)$, $f(t^{2})$ and $f(t-2)$, and state the degree of each in $t$. \bk{Ex 1.2 Q10}
\item The number of handshakes among $x$ people is $H(x)=\dfrac x2(x-1)$. Find $H(5)$, $H(6)$, $H(10)$, and the number of people who give 136 handshakes. \bk{Ex 1.2 Q15}
\end{qlist}

{\color{lilacdeep}\footnotesize\bfseries FACTORISING (1.3)}
\begin{qlist}[resume]
\item Factorise $3ax^{2}-3ay^{2}-4bx^{2}+4by^{2}$. \bk{Ex 1.3 Q17}
\item Factorise $12x^{2}+17xy-5y^{2}$. \bk{Ex 1.3 Q50}
\item Using the quadratic formula, factorise $2x^{2}-5\sqrt2x-6$. \bk{Ex 1.3 Q51}
\item Factorise $(x+y)^{3}-z^{3}$. \bk{Ex 1.3 Q55}
\end{qlist}

{\color{lilacdeep}\footnotesize\bfseries ALGEBRAIC FRACTIONS (1.4)}
\begin{qlist}[resume]
\item Simplify $\dfrac{t^{2}+3t-4}{t^{2}-3t+2}$. \bk{Ex 1.4 Q3}
\item Simplify $\dfrac{1}{x^{2}-9}-\dfrac{2}{x^{2}-x-6}$. \bk{Ex 1.4 Q5}
\item Show that $\dfrac{3x-5}{x-2}+\dfrac{1}{2-x}$ simplifies to a constant when $x\neq2$. \bk{Ex 1.4 Q12}
\end{qlist}

{\color{lilacdeep}\footnotesize\bfseries BINOMIAL EXPANSIONS (1.5)}
\begin{qlist}[resume]
\item If $\binom{12}{k}=\binom{12}{3}$, find $k\in\mathbb N$, $k\neq3$. \bk{Ex 1.5 Q3}
\item How many terms are in the expansion of $(p+2q)^{6}$? Hence find the coefficient of the middle term. \bk{Ex 1.5 Q14}
\item Find the coefficient of the 4th term of $\left(x-\dfrac{3y}{2}\right)^{9}$. \bk{Ex 1.5 Q17}
\end{qlist}

{\color{lilacdeep}\footnotesize\bfseries IDENTITIES (1.6)}
\begin{qlist}[resume]
\item If $x^{2}-4x-5=(x-n)^{2}-m$ for all $x$, find $m$ and $n$. \bk{Ex 1.6 Q7}
\item Given that $(x^{2}-4)$ is a factor of $x^{3}+cx^{2}+dx-12$, find $c$ and $d$, and factorise the cubic fully. \bk{Ex 1.6 Q22}
\item Write $\dfrac{1}{(x+1)(x+4)}$ as partial fractions. \bk{Ex 1.6 Q19}
\end{qlist}

{\color{lilacdeep}\footnotesize\bfseries MANIPULATING FORMULAE (1.7)}
\begin{qlist}[resume]
\item A cylinder has $V=\pi r^{2}h$ and curved surface area $A=2\pi rh$. Find $r$ in terms of $V,h$ and in terms of $A,h$; hence show $A^{2}=4\pi hV$. \bk{Ex 1.7 Q3}
\item The value of \texteuro$P$ after 3 years at $i\%$ is $A=P\left(1+\dfrac{i}{100}\right)^{3}$. Find $i$ in terms of $P$ and $A$; if \texteuro2500 is now worth \texteuro2650, find $i$ to one decimal place. \bk{Ex 1.7 Q9}
\item A farmer uses 300\,m of fencing for a rectangular paddock against a wall. Find $L$ in terms of $W$, then the area $A$ in terms of $W$, and the dimensions giving $A=10\,000$\,m$^{2}$. \bk{Ex 1.7 Q12}
\end{qlist}

{\color{lilacdeep}\footnotesize\bfseries PATTERNS (1.8)}
\begin{qlist}[resume]
\item Find the rule for the pattern $6,12,20,30,42,\dots$ using first and second differences, and hence the 100th term. \bk{Rev.\ Core Q13}
\item Bacteria counted hourly give $4,7,14,25,40,\dots$. Show the pattern has both a linear and a quadratic element, find $f(t)$, and find the hour in which the colony reaches 529. \bk{Ex 1.8 Q8}
\end{qlist}

{\color{lilacdeep}\footnotesize\bfseries SOLVING EQUATIONS (1.9)}
\begin{qlist}[resume]
\item A game sells for \texteuro11.50; production costs \texteuro3500 plus \texteuro10.50 per game. Write $C(x)$ and $I(x)$, find how many games recoup the costs, and how many give a profit of \texteuro2000. \bk{Ex 1.9 Q4}
\end{qlist}

{\color{lilacdeep}\footnotesize\bfseries SIMULTANEOUS EQUATIONS (1.10)}
\begin{qlist}[resume]
\item Solve $2x+y+z=8$, $5x-3y+2z=3$, $7x+y+3z=20$. \bk{Ex 1.10 Q5}
\item The curve $f(x)=ax^{2}+bx+c$ passes through $(1,2)$, $(2,4)$ and $(3,8)$. Find $f(x)$. \bk{Ex 1.10 Q8}
\item How many litres of 70\% alcohol must be added to 50 litres of 40\% alcohol to make a 50\% solution? \bk{Ex 1.10 Q16}
\item The circle $x^{2}+y^{2}+ax+by+c=0$ passes through $(1,0)$, $(1,2)$ and $(2,1)$. Find $a$, $b$ and $c$. \bk{Ex 1.10 Q22}
\end{qlist}

{\color{lilacdeep}\footnotesize\bfseries REVISION}
\begin{qlist}[resume]
\item If $p(x-q)^{2}+r=2x^{2}-12x+5$ for all $x$, find $p$, $q$ and $r$. \bk{Rev.\ Core Q9}
\item A closed box has a square base of side $x$\,cm, height $h$\,cm and volume 40\,cm$^{3}$. Write $h$ in terms of $x$ and show that the surface area is $S=2x^{2}+\dfrac{160}{x}$. \bk{Rev.\ Ext.\ Q3}
\end{qlist}
\end{multicols}

\clearpage

%%%%%%%%%%%%%%%%%%%%%%%%%%%%%%%  APPENDIX: SOLUTIONS  %%%%%%%%%%%%%%%%%%%%%%%%%%%%%%%

{\color{lilacdeep}\large\bfseries Appendix --- Solutions}\hfill
{\color{greyish}\footnotesize Q1--Q33, with the working expected in an answer}\\[-5pt]

\textcolor{lilacline}{\rule{\textwidth}{0.9pt}}

\solhead{1}{polynomial or not}
$\dfrac4x=4x^{-1}$ has a \emph{negative} power, and $x^{\frac32}$ has a \emph{fractional} power.
\key{A polynomial may contain only whole-number powers of the variable}, so either reason alone disqualifies it.

\solhead{2}{perfect square}
$25x^{2}=(5x)^{2}$ and $4=(\pm2)^{2}$, so the square is $(5x\pm2)^{2}=25x^{2}\pm20x+4$.
Comparing middle terms, $tx=\pm20x$, so $t=20$ or $t=-20$. (Equivalently $t=\pm2\sqrt{ac}=\pm2\sqrt{25\cdot4}$.)

\solhead{3}{picking out one coefficient}
Do not expand everything --- collect only the products that give $x^{2}$:
\[(3x^{2})(-5)+(5x)(-6x)+(-1)(2x^{2})=-15x^{2}-30x^{2}-2x^{2}=-47x^{2},\]
so the coefficient is $\mathbf{-47}$.

\solhead{4}{long division}
\vspace{-0.6em}
\begin{center}
\begin{tabular}{@{}r@{\ }l@{\qquad}l@{}}
& \underline{$x^{2}-7x+12$} & \\
$x-1\,\big)$ & $x^{3}-8x^{2}+19x-12$ & \small\color{greyish}$x^{3}\div x=x^{2}$\\
& \underline{$x^{3}-\ x^{2}$}\hspace{2.75cm} & \small\color{greyish}multiply, subtract\\
& $\ \ \ \ -7x^{2}+19x-12$ & \small\color{greyish}$-7x^{2}\div x=-7x$\\
& $\ \ \ \ \underline{-7x^{2}+7x}$\hspace{1.5cm} & \small\color{greyish}multiply, subtract\\
& $\ \ \ \ \ \ \ \ \ \ \ \ \ \ 12x-12$ & \small\color{greyish}$12x\div x=12$\\
& $\ \ \ \ \ \ \ \ \ \ \ \ \ \ \underline{12x-12}$ & \small\color{greyish}remainder $0$\\
\end{tabular}
\end{center}
\vspace{-0.4em}
So the quotient is $x^{2}-7x+12$, and since the remainder is zero,
$x^{3}-8x^{2}+19x-12=(x-1)(x-3)(x-4)$.

\solhead{5}{difference of two cubes}
$8x^{3}-27y^{3}=(2x)^{3}-(3y)^{3}=(2x-3y)\!\left[(2x)^{2}+(2x)(3y)+(3y)^{2}\right]$, so
\[\frac{8x^{3}-27y^{3}}{2x-3y}=4x^{2}+6xy+9y^{2}.\]
\key{Recognising the cubes is quicker than dividing}, though long division gives the same answer.

\solhead{6}{evaluating a function}
$d(n)=\frac{n^{2}}{2}-\frac{3n}{2}=\frac{n(n-3)}{2}$, so
$d(4)=\frac{4(1)}{2}=2$, \ $d(5)=\frac{5(2)}{2}=5$, \ $d(6)=\frac{6(3)}{2}=9$.
$d(3)=\frac{3(0)}{2}=0$: in a triangle every pair of vertices is already joined by a side, so there is no diagonal to draw.

\solhead{7}{substituting expressions}
$f(-2t)=(-2t)^{2}-3(-2t)+6=4t^{2}+6t+6$ \ (degree 2);\\
$f(t^{2})=(t^{2})^{2}-3t^{2}+6=t^{4}-3t^{2}+6$ \ (degree 4);\\
$f(t-2)=(t-2)^{2}-3(t-2)+6=t^{2}-4t+4-3t+6+6=t^{2}-7t+16$ \ (degree 2).\\
\key{Keep the substituted expression in brackets} --- that is where the sign errors happen.

\solhead{8}{a quadratic model}
$H(5)=\frac52(4)=10$, \ $H(6)=\frac62(5)=15$, \ $H(10)=\frac{10}{2}(9)=45$.\\
For 136 handshakes: $\frac x2(x-1)=136\Rightarrow x^{2}-x-272=0\Rightarrow(x-17)(x+16)=0$, so $x=17$ (the root $x=-16$ is rejected --- a count of people must be positive). \textbf{17 students.}

\solhead{9}{grouping, then difference of two squares}
\[3ax^{2}-3ay^{2}-4bx^{2}+4by^{2}=3a(x^{2}-y^{2})-4b(x^{2}-y^{2})=(x^{2}-y^{2})(3a-4b)=(x-y)(x+y)(3a-4b).\]

\solhead{10}{quadratic trinomial in two letters}
Treat it as a quadratic in $x$: the factor pairs of $12$ and $-5$ that give a middle term $+17xy$ are
$12x^{2}+17xy-5y^{2}=(3x+5y)(4x-y)$.
Check the middle term: $-3xy+20xy=+17xy$.\ \checkmark

\solhead{11}{formula with surds}
$a=2$, $b=-5\sqrt2$, $c=-6$, so $b^{2}-4ac=50+48=98$ and $\sqrt{98}=7\sqrt2$:
\[x=\frac{5\sqrt2\pm7\sqrt2}{4}=\frac{12\sqrt2}{4}\ \text{or}\ \frac{-2\sqrt2}{4}=3\sqrt2\ \text{or}\ -\frac{\sqrt2}{2}.\]
The root $3\sqrt2$ gives the factor $(x-3\sqrt2)$; the root $-\frac{\sqrt2}{2}$ gives $(2x+\sqrt2)$. Hence
$2x^{2}-5\sqrt2x-6=(x-3\sqrt2)(2x+\sqrt2)$.
\key{The leading coefficient 2 has to go somewhere} --- put it with the fractional root.

\solhead{12}{difference of two cubes with a bracket}
Use the template with ``$x$''$=(x+y)$ and ``$y$''$=z$:
\[(x+y)^{3}-z^{3}=\big[(x+y)-z\big]\big[(x+y)^{2}+(x+y)z+z^{2}\big]=(x+y-z)(x^{2}+2xy+y^{2}+xz+yz+z^{2}).\]

\solhead{13}{cancel factors}
\[\frac{t^{2}+3t-4}{t^{2}-3t+2}=\frac{(t+4)(t-1)}{(t-1)(t-2)}=\frac{t+4}{t-2}.\]

\solhead{14}{factorise before choosing the denominator}
$x^{2}-9=(x-3)(x+3)$ and $x^{2}-x-6=(x-3)(x+2)$, so the common denominator is $(x-3)(x+3)(x+2)$:
\[\frac{(x+2)-2(x+3)}{(x-3)(x+3)(x+2)}=\frac{x+2-2x-6}{(x-3)(x+3)(x+2)}=\frac{-(x+4)}{(x-3)(x+3)(x+2)}.\]

\solhead{15}{a sign trick}
\key{Note that \mbox{$2-x=-(x-2)$}}, so $\dfrac{1}{2-x}=-\dfrac{1}{x-2}$. Then
\[\frac{3x-5}{x-2}-\frac{1}{x-2}=\frac{3x-6}{x-2}=\frac{3(x-2)}{x-2}=3,\]
a constant for every $x\neq2$.

\solhead{16}{symmetry of the binomial coefficients}
$\binom nr=\binom n{n-r}$, so $\binom{12}{3}=\binom{12}{9}$ and $k=\mathbf{9}$
(both equal 220).

\solhead{17}{the middle term}
$(p+2q)^{6}$ has $6+1=\mathbf7$ terms, so the middle one is the 4th, i.e.\ $r=3$:
\[\binom63p^{3}(2q)^{3}=20\cdot8\,p^{3}q^{3}=160p^{3}q^{3},\]
and the coefficient is $\mathbf{160}$.

\solhead{18}{a named term}
The 4th term means $r=3$ in $\binom nrx^{n-r}y^{r}$, with $n=9$ and $y=-\frac{3y}{2}$:
\[\binom93x^{6}\left(-\frac{3y}{2}\right)^{3}=84\,x^{6}\left(-\frac{27}{8}y^{3}\right)=-\frac{567}{2}x^{6}y^{3},\]
so the coefficient is $-\frac{567}{2}=-283.5$. \key{The cube of a negative term stays negative}, and the $2$ in the denominator is cubed too.

\solhead{19}{completing the square as an identity}
$(x-n)^{2}-m=x^{2}-2nx+n^{2}-m$. Comparing with $x^{2}-4x-5$ for all $x$:
\[-2n=-4\Rightarrow n=2,\qquad n^{2}-m=-5\Rightarrow4-m=-5\Rightarrow m=9.\]
Check: $(x-2)^{2}-9=x^{2}-4x+4-9=x^{2}-4x-5$.\ \checkmark

\solhead{20}{a known factor}
If $(x^{2}-4)=(x-2)(x+2)$ is a factor, then $x=2$ and $x=-2$ both give zero:
\[x=2:\ 8+4c+2d-12=0\Rightarrow2c+d=2;\qquad x=-2:\ -8+4c-2d-12=0\Rightarrow2c-d=10.\]
Adding, $4c=12\Rightarrow c=3$, and then $d=2-2(3)=-4$. The cubic is $x^{3}+3x^{2}-4x-12$, and dividing by $x^{2}-4$ gives $x+3$, so
\[x^{3}+3x^{2}-4x-12=(x-2)(x+2)(x+3).\]

\solhead{21}{partial fractions}
$\dfrac{1}{(x+1)(x+4)}=\dfrac{A}{x+1}+\dfrac{B}{x+4}\Rightarrow1=A(x+4)+B(x+1)$ for all $x$.
Substituting convenient values: $x=-1\Rightarrow1=3A\Rightarrow A=\frac13$; \ $x=-4\Rightarrow1=-3B\Rightarrow B=-\frac13$. Hence
\[\frac{1}{(x+1)(x+4)}=\frac{1}{3(x+1)}-\frac{1}{3(x+4)}.\]

\solhead{22}{two rearrangements and a link between them}
From $V=\pi r^{2}h$: $r^{2}=\dfrac{V}{\pi h}\Rightarrow r=\sqrt{\dfrac{V}{\pi h}}$. \ From $A=2\pi rh$: $r=\dfrac{A}{2\pi h}$.
Equating the two and squaring, $\dfrac{A^{2}}{4\pi^{2}h^{2}}=\dfrac{V}{\pi h}$, so
\[A^{2}=4\pi^{2}h^{2}\cdot\frac{V}{\pi h}=4\pi hV. \qquad\checkmark\]

\solhead{23}{a root and a percentage}
$\dfrac AP=\left(1+\dfrac{i}{100}\right)^{3}\Rightarrow\sqrt[3]{\dfrac AP}=1+\dfrac{i}{100}
\Rightarrow i=100\left(\sqrt[3]{\dfrac AP}-1\right).$
With $A=2650$ and $P=2500$: $\frac AP=1.06$, $\sqrt[3]{1.06}=1.01961$, so $i=1.961\approx\mathbf{2.0\%}$.

\solhead{24}{modelling, then a quadratic}
The wall replaces one length, so the fencing is $L+2W=300$, giving $L=300-2W$. Then
\[A=LW=W(300-2W)=300W-2W^{2}.\]
Setting $A=10\,000$: $2W^{2}-300W+10\,000=0\Rightarrow W^{2}-150W+5000=0\Rightarrow(W-50)(W-100)=0$.
So $W=50$ with $L=200$, or $W=100$ with $L=100$ --- \key{two different paddocks of the same area}.

\solhead{25}{second differences}
\vspace{-0.5em}
\begin{center}
\begin{tabular}{@{}l|ccccc@{}}
\textbf{pattern} & 6 & 12 & 20 & 30 & 42\\
\textbf{1st difference} & & 6 & 8 & 10 & 12\\
\textbf{2nd difference} & & & 2 & 2 & 2\\
\end{tabular}
\end{center}
\vspace{-0.3em}
The 2nd difference is constant, so the pattern is quadratic, $f(x)=ax^{2}+bx+c$ with $x\ge0$:
$2a=2\Rightarrow a=1$; \ $c=f(0)=6$; \ $f(1)=1+b+6=12\Rightarrow b=5$. So
\[f(x)=x^{2}+5x+6=(x+2)(x+3).\]
\key{The pattern is indexed from \mbox{$x=0$}, so the 100th term is \mbox{$f(99)$}}$=9801+495+6=\mathbf{10\,302}$.

\solhead{26}{a pattern with both elements}
1st differences $3,7,11,15$ (not constant --- so there is a linear part changing); 2nd differences $4,4,4$ (constant --- so there is a quadratic part). With $f(t)=at^{2}+bt+c$:
$2a=4\Rightarrow a=2$; \ $c=f(0)=4$; \ $f(1)=2+b+4=7\Rightarrow b=1$. So
\[f(t)=2t^{2}+t+4,\]
which reproduces $14$, $25$ and $40$ at $t=2,3,4$. \checkmark\ \ Trying values: $f(15)=450+15+4=469$ and $f(16)=512+16+4=532$, so the colony passes 529 \key{during the 16th hour}.

\solhead{27}{cost, income and profit}
$C(x)=3500+10.5x$ and $I(x)=11.5x$. The costs are recouped when the graphs cross:
\[11.5x=3500+10.5x\Rightarrow x=\mathbf{3500}\ \text{games}.\]
$P=I-C=11.5x-(3500+10.5x)=x-3500$, which is the profit. For $P=2000$: $x-3500=2000\Rightarrow x=\mathbf{5500}$ games.
(Each game contributes only \texteuro1 of margin, which is why the numbers are so large.)

\solhead{28}{three unknowns}
A: $2x+y+z=8$; \ B: $5x-3y+2z=3$; \ C: $7x+y+3z=20$. Eliminate $z$ twice:
\[2\text{A}-\text{B}:\ -x+5y=13,\qquad 3\text{A}-\text{C}:\ -x+2y=4.\]
Subtracting, $3y=9\Rightarrow y=3$; then $-x+2(3)=4\Rightarrow x=2$; then A gives $4+3+z=8\Rightarrow z=1$.
\[(x,y,z)=(2,3,1),\]
which satisfies all three: $4+3+1=8$, $10-9+2=3$, $14+3+3=20$.\ \checkmark

\solhead{29}{fitting a curve through three points}
Each point gives an equation in $a$, $b$, $c$:
\[a+b+c=2,\qquad 4a+2b+c=4,\qquad 9a+3b+c=8.\]
Subtracting in pairs: $3a+b=2$ and $5a+b=4$, so $2a=2\Rightarrow a=1$, $b=-1$, $c=2$. Hence
$f(x)=x^{2}-x+2$, and indeed $f(3)=9-3+2=8$.\ \checkmark

\solhead{30}{a mixture}
Let $x$ litres of the 70\% solution be added. Counting the \emph{alcohol}, not the liquid:
\[0.7x+0.4(50)=0.5(x+50)\Rightarrow0.7x+20=0.5x+25\Rightarrow0.2x=5\Rightarrow x=\mathbf{25}\ \text{litres}.\]
\key{The equation always balances the quantity of the thing being mixed}, here alcohol.

\solhead{31}{a circle through three points}
Substituting each point into $x^{2}+y^{2}+ax+by+c=0$:
\[(1,0):\ 1+a+c=0;\qquad(1,2):\ 5+a+2b+c=0;\qquad(2,1):\ 5+2a+b+c=0.\]
From the first, $a+c=-1$. Putting that into the second: $-1+2b=-5\Rightarrow b=-2$. The third gives $2a+c=-3$; subtracting $a+c=-1$ leaves $a=-2$, so $c=1$.
\[a=-2,\quad b=-2,\quad c=1\qquad\text{i.e.}\qquad x^{2}+y^{2}-2x-2y+1=0.\]

\solhead{32}{an identity in disguise}
$p(x-q)^{2}+r=px^{2}-2pqx+pq^{2}+r$. Comparing with $2x^{2}-12x+5$ for all $x$:
\[p=2;\qquad-2pq=-12\Rightarrow-4q=-12\Rightarrow q=3;\qquad pq^{2}+r=5\Rightarrow18+r=5\Rightarrow r=-13.\]

\solhead{33}{eliminating a variable from a model}
The volume fixes $h$: $x^{2}h=40\Rightarrow h=\dfrac{40}{x^{2}}$. A closed box has two square ends and four rectangular sides:
\[S=2x^{2}+4xh=2x^{2}+4x\left(\frac{40}{x^{2}}\right)=2x^{2}+\frac{160}{x}. \qquad\checkmark\]
\key{Use the fixed quantity to remove one variable}, leaving a function of $x$ alone --- exactly the step that sets up an optimisation question later in the course.

\end{document}
